\item \textbf{Sistema de aterrizaje por instrumentos con doble antena.}
Su grupo está estudiando el sistema de aterrizaje por instrumentos utilizado para guiar a los aviones a aterrizajes seguros en algunos aeropuertos cuando la visibilidad es pobre. Resulta que el experimento de doble rendija de Young se usa en este sistema. Un piloto intenta alinear su avión con una pista. Dos antenas de radio (los puntos negros en la figura) están colocadas adyacentes a la pista, separadas una distancia $d = 40.0\text{ m}$. Las antenas transmiten ondas de radio coherentes no moduladas a $30.0\text{ MHz}$. Las líneas en la figura TP36.1 representan trayectorias a lo largo de las cuales existe un patrón de interferencia máxima.
\begin{enumerate}
    \item Encuentre la longitud de onda de las ondas de radio.
    \item \textbf{¿Qué pasaría si?} Suponga que el avión vuela a lo largo del primer máximo lateral en vez del máximo central. ¿A qué distancia de la lateral de la línea de centro de la pista estará el avión cuando esté a $2.00\text{ km}$ de las antenas medidos a lo largo de su dirección de vuelo?
    \item Es posible avisarle al piloto que está en el máximo equivocado si se le envían dos señales desde cada antena y se equipa al avión con un receptor de dos canales. La razón entre las dos frecuencias no debe ser la razón entre enteros pequeños (por ejemplo, $3/4$). Explique la forma en que funcionaría este sistema de dos frecuencias y por qué no daría resultado si las frecuencias estuvieran relacionadas por una razón de enteros sencillos.
\end{enumerate}